% Part: first-order-logic
% Chapter: model-theory
% Section: isomorphism

\documentclass[../../../include/open-logic-section]{subfiles}

\begin{document}

\olfileid{mod}{bas}{iso}
\olsection{Dezelfde structuur met andere objecten (isomorfe structuren)}

!!{structure}s van de eerste orde kunnen op twee verschillende manieren
hetzelfde lijken. Bij de eerste manier zijn in beide precies dezelfde
!!{sentence}s waar. Zulke !!{structure}s heten \emph{elementair
equivalent}. Toch kunnen structuren dan nog heel verschillend zijn. De
ene kan bijvoorbeeld !!{enumerable} zijn en de andere niet, terwijl ze
dezelfde !!{sentence}s waar maken. Veel eigenschappen van
!!a{structure} kun je namelijk niet in een taal van de eerste orde
uitdrukken. Soms heeft de taal daarvoor niet genoeg middelen. Soms komt
dat door een grens van de eerste-ordelogica zelf, zoals de stelling van
L\"owenheim--Skolem. Er is daarom een tweede, strengere manier waarop
!!{structure}s hetzelfde kunnen zijn. Dan is hun hele structurele
patroon gelijk en zijn alleen de concrete objecten anders. Een
\emph{isomorfisme} legt precies vast wat dat betekent.

\begin{defn}
\ollabel{defn:elem-equiv} 
Neem twee !!{structure}s $\Struct{M}$ en $\Struct M'$ voor dezelfde
!!{language}~$\Lang{L}$. We zeggen dat $\Struct{M}$ \emph{elementair
equivalent} is met $\Struct M'$, geschreven als $\Struct{M} \equiv
\Struct M'$, als in beide precies dezelfde zinnen waar zijn. Formeel
betekent dit: voor elke !!{sentence}~$!A$ van~$\Lang{L}$,
$\Sat{M}{!A}$ dan en slechts dan als $\Sat{M'}{!A}$.
\end{defn}

\begin{defn}
\ollabel{defn:isomorphism} Neem twee !!{structure}s $\Struct{M}$ en
$\Struct M'$ voor dezelfde !!{language}~$\Lang L$. We zeggen dat
$\Struct{M}$ \emph{isomorf} is met $\Struct M'$, oftewel dat beide
structureel hetzelfde zijn. We schrijven $\Struct{M}
\simeq \Struct M'$ als er een functie $h \colon
\Domain{M} \to \Domain{M'}$ bestaat met de volgende eigenschappen:
\begin{enumerate}
\item $h$ is !!{injective}: twee verschillende objecten krijgen nooit
  hetzelfde beeld; dus als $h(x) = h(y)$, dan $x = y$; 
\item $h$ is !!{surjective}: elk object aan de andere kant wordt geraakt;
  voor elke $y \in \Domain{M'}$ is er dus een $x \in \Domain{M}$ waarvoor
  $h(x) = y$;
\item \ollabel{defn:iso-const}voor elke !!{constant} $c$ wordt het ene
  aangewezen object op het andere afgebeeld:
  $h(\Assign{c}{M}) = \Assign{c}{M'}$;
\item \ollabel{defn:iso-pred}voor elk $n$-plaatsig !!{predicate}~$P$
  geldt in beide richtingen dat een rij objecten onder het predicaat valt:
  \[
  \tuple{a_1, \dots, a_n}\in \Assign{P}{M} \quad\text{iff}\quad
  \tuple{h(a_1), \dots, h(a_n)} \in \Assign{P}{M'};
  \]
\item \ollabel{defn:iso-func}voor elke $n$-plaatsige !!{function} $f$
  maakt het niet uit of je eerst de functie in de eerste structuur
  gebruikt en de uitkomst omzet, of eerst alle invoeren omzet en daarna
  de functie in de tweede structuur gebruikt:
  \[
  h(\Assign{f}{M}(a_1, \dots, a_n)) =
  \Assign{f}{M'}(h(a_1), \dots, h(a_n)).
  \]
\end{enumerate}
\end{defn}

\begin{thm}
\ollabel{thm:isom}
Als $\Struct{M} \iso \Struct M'$, dan $\Struct{M} \elemequiv
\Struct{M'}$.
\end{thm}

\begin{proof}
Neem een isomorfisme $h$ van $\Struct{M}$ op $\Struct M'$. Een
toekenning~$s$ wijst aan elke variabele een object toe.
De samenstelling $h \circ s$ combineert $h$ en $s$: eerst wordt de
toekenning gebruikt en daarna het isomorfisme. Het is dus de toekenning
in $\Struct{M'}$ waarvoor
$(h \circ s)(x) = h(s(x))$. Met inductie naar de opbouw van $t$ en $!A$
kunnen we tegelijk de volgende twee, sterkere uitspraken bewijzen:
\begin{enumerate}
  \item[a.] $h(\Value{t}{M}[s]) = \Value{t}{M'}[h\circ s]$.
  \item[b.] $\Sat{M}{!A}[s]$ dan en slechts dan als $\Sat{M'}{!A}[h \circ s]$.
\end{enumerate}
We bewijzen de eerste uitspraak met inductie naar de opbouw van~$t$.
\begin{enumerate}
\item Als $t \ident c$, dan $\Value{c}{M}[s] = \Assign{c}{M}$ en
  $\Value{c}{M'}[h \circ s] = \Assign{c}{M'}$. Daarom geldt
  $h(\Value{t}{M}[s]) = h(\Assign{c}{M}) = \Assign{c}{M'}$ (volgens
  \olref{defn:iso-const} van \olref{defn:isomorphism}) $=
  \Value{t}{M'}[h \circ s]$.
\item Als $t \ident x$, dan $\Value{x}{M}[s] = s(x)$ en
  $\Value{x}{M'}[h \circ s] = h(s(x))$. Daarom is $h(\Value{x}{M}[s]) =
  h(s(x)) = \Value{x}{M'}[h \circ s]$.
\item Als $t \ident f(t_1, \dots, t_n)$, dan
  \begin{align*}
    \Value{t}{M}[s] & = \Assign{f}{M}(\Value{t_1}{M}[s], \dots, \Value{t_n}{M}[s]) \quad\text{and}\\
  \Value{t}{M'}[h \circ s] & = \Assign{f}{M}(\Value{t_1}{M'}[h \circ
    s], \dots, \Value{t_n}{M'}[h \circ s]).
  \end{align*}
  De inductiehypothese zegt voor elke $i$ dat
  $h(\Value{t_i}{M}[s]) = \Value{t_i}{M'}[h\circ s]$. Daarmee krijgen
  we:
  \begin{align}
    h(\Value{t}{M}[s]) 
    & = h(\Assign{f}{M}(\Value{t_1}{M}[s], \dots, \Value{t_n}{M}[s]) \notag\\
    & = \Assign{f}{M'}(h(\Value{t_1}{M}[s]), \dots,
    h(\Value{t_n}{M}[s])) \ollabel{iso-1}\\
    & = \Assign{f}{M'}(\Value{t_1}{M'}[h \circ s], \dots,
    \Value{t_n}{M'}[h \circ s]) \ollabel{iso-2}\\
    & = \Value{t}{M'}[h\circ s] \notag
  \end{align}
  De stap \olref{iso-1} gebruikt \olref{defn:iso-func} uit
  \olref{defn:isomorphism}. De stap \olref{iso-2} gebruikt de
  inductiehypothese.
\end{enumerate}
Het bewijs van (b) is een oefening.

Als $!A$ een zin is, maakt de keuze van de toekenningen~$s$ en
$h \circ s$ niet uit. Daarom geldt $\Sat{M}{!A}$ dan en slechts dan als
$\Sat{M'}{!A}$.
\end{proof}

\begin{prob}
Werk het bewijs van (b) van \olref[mod][bas][iso]{thm:isom} stap voor
stap uit. Geef aan waar je elk van de vijf eigenschappen uit
\olref[mod][bas][iso]{defn:isomorphism} gebruikt die samen bepalen wat
een isomorfisme is.
\end{prob}

\begin{defn}
Een \emph{automorfisme} van een structuur $\Struct{M}$ is een
isomorfisme van $\Struct{M}$ naar zichzelf.
\end{defn}

\begin{prob}
Toon voor elke structuur $\Struct{M}$ het volgende aan. Stel dat $X$ een
deelverzameling van $\Struct{M}$ is die met een formule is vastgelegd,
dus een definieerbare deelverzameling, en dat $h$ een automorfisme van
$\Struct{M}$ is. Dan $X = \Setabs{h(x)}{x \in X}$ (dat wil zeggen: $X$
blijft dezelfde verzameling als we $h$ op alle elementen toepassen).
\end{prob}


\end{document}
